\documentclass[11pt]{article}
\usepackage[margin=0.92in]{geometry}
\usepackage{amsmath,amssymb,amsthm,mathtools}
\usepackage{graphicx}
\usepackage{xcolor}
\usepackage{microtype}
\usepackage[hidelinks]{hyperref}
\emergencystretch=3em

\newtheorem{theorem}{Theorem}
\newtheorem{lemma}[theorem]{Lemma}
\newtheorem{remark}[theorem]{Remark}
\newcommand{\D}{\mathbb D}
\newcommand{\T}{\mathbb T}

\title{Boundary-regular isometric symbols on analytic Besov spaces}
\author{Candidate partial result for arXiv:2207.12634}
\date{17 August 2026}

\begin{document}
\maketitle

\begin{center}
\fbox{\parbox{0.91\textwidth}{\textbf{Status.}
Candidate partial result, likely valid.  The unrestricted conjecture for an
arbitrary analytic symbol remains open in this packet.  We prove it when the
symbol extends holomorphically past the closed unit disk.}}
\end{center}

\section{Source conjecture}

For $1<p<\infty$, the analytic Besov space $B_p$ consists of the analytic
functions on $\D$ for which
\[
 \|f\|_p^p=\int_\D |f'(z)|^p(1-|z|^2)^{p-2}\,dA(z)<\infty,
\]
with Banach norm $|f(0)|+\|f\|_p$.  Allen, Heller, and Pons
\cite{AHP} proved that rotations are the only isometric composition symbols
for $1<p<2$, and obtained several reductions for $p>2$.  They ask whether the
same conclusion holds without an extra hypothesis when $p>2$.

\begin{center}
\includegraphics[width=0.96\textwidth]{figures/open_problem_crop.png}
\end{center}

The theorem below answers the question for symbols with analytic boundary
regularity.  In particular it covers rational self-maps with no pole on
$\overline\D$.

\section{Idea of the proof}

The isometry condition supplies two rigid pieces of information: the symbol
is full up to area zero, and the hyperbolic derivative weights over every
fiber sum to one.  Analytic continuation past the circle turns fullness into
the boundary identity $|\varphi|=1$, so the symbol is a finite Blaschke
product.  Near a boundary value, every inverse branch of such a product has
hyperbolic derivative tending to one.  The fiber identity then says that the
degree of the product is one.  The remaining normalization
$\varphi(0)=0$ makes the resulting automorphism a rotation.

\section{The fiber identity}

Fix $p>2$ and write $\beta=p-2>0$.  If $C_\varphi$ is an isometry on $B_p$,
the source proves that
\begin{equation}\label{eq:zero-full}
 \varphi(0)=0,
 \qquad A(\D\setminus\varphi(\D))=0,
\end{equation}
and that $W_{\varphi',\varphi}$ is an isometry on the weighted Bergman space
$A^p_\beta$; see \cite[Lemmas 3.1 and 3.4]{AHP}.  Kolaski's Borel
equimeasurability identity, followed by the area formula, therefore gives
\begin{equation}\label{eq:fiber}
 \sum_{\varphi(z)=w}
 \left(
 \frac{|\varphi'(z)|(1-|z|^2)}{1-|w|^2}
 \right)^\beta=1
\end{equation}
for almost every regular value $w\in\D$.  Each summand is at most one by the
Schwarz--Pick lemma.  Formula \eqref{eq:fiber} is the only isometric input
needed after \eqref{eq:zero-full}.

\section{Full boundary-regular maps}

\begin{lemma}\label{lem:inner}
Let $\varphi:\D\to\D$ extend holomorphically to a neighborhood of
$\overline\D$.  If $A(\D\setminus\varphi(\D))=0$, then
$|\varphi(\zeta)|=1$ for every $\zeta\in\T$.  Consequently $\varphi$ is a
finite Blaschke product up to a unimodular constant.
\end{lemma}

\begin{proof}
Polar Fubini applied to the null set $\D\setminus\varphi(\D)$ shows that for
almost every $\eta\in\T$ there are radii $r_n\uparrow1$ with
$r_n\eta\in\varphi(\D)$.  Choose $z_n\in\D$ with
$\varphi(z_n)=r_n\eta$.  A subsequence converges in $\overline\D$ to some
$\zeta$.  Continuity gives $\varphi(\zeta)=\eta$.  The point $\zeta$ cannot
lie in $\D$, by the maximum-modulus principle, so $\zeta\in\T$.

Thus the image under $\varphi|_\T$ of
\[
 Z=\{\zeta\in\T:|\varphi(\zeta)|=1\}
\]
contains almost all of $\T$.  The boundary map is Lipschitz.  Hence $Z$
cannot have arc-length measure zero, since Lipschitz maps send null sets to
null sets.  The function
$t\mapsto |\varphi(e^{it})|^2-1$ is real analytic and vanishes on the
positive-measure set corresponding to $Z$; it therefore vanishes identically.
The standard zero-factorization and maximum principle now show that
$\varphi$ is a finite Blaschke product times a unimodular constant.
\end{proof}

\section{Rotation theorem}

\begin{theorem}\label{thm:main}
Let $p>2$, let $\varphi:\D\to\D$ extend holomorphically to a neighborhood of
$\overline\D$, and suppose $C_\varphi$ is an isometry on $B_p$.  Then there is
$\theta\in\mathbb R$ such that
\[
 \varphi(z)=e^{i\theta}z\qquad(z\in\D).
\]
\end{theorem}

\begin{proof}
By \eqref{eq:zero-full} and Lemma~\ref{lem:inner}, $\varphi$ is a finite
Blaschke product; let its degree be $d\ge1$.  Choose a point $\eta\in\T$ and
a sequence $w_n\to\eta$ such that every $w_n$ is a regular value and
\eqref{eq:fiber} holds.  This is possible because the set on which
\eqref{eq:fiber} holds has full area and is therefore dense, and a finite
Blaschke product has only finitely many critical values.

For large $n$, the equation $\varphi(z)=w_n$ has $d$ simple roots
$z_{1,n},\ldots,z_{d,n}$.  After labeling and passing to subsequences, each
$z_{j,n}$ tends to a boundary point $\zeta_j$ with
$\varphi(\zeta_j)=\eta$.  A finite Blaschke product has a finite nonzero
angular derivative at every point of $\T$.  The Julia--Carath\'eodory
identity, applied along the inverse branch, gives
\[
 \frac{|\varphi'(z_{j,n})|(1-|z_{j,n}|^2)}{1-|w_n|^2}
 \longrightarrow 1
 \qquad (j=1,\ldots,d).
\]
Passing to the limit in \eqref{eq:fiber} yields
\[
 1=\lim_{n\to\infty}\sum_{j=1}^d
 \left(
 \frac{|\varphi'(z_{j,n})|(1-|z_{j,n}|^2)}{1-|w_n|^2}
 \right)^\beta=d.
\]
Thus $d=1$, so $\varphi$ is a disk automorphism.  Finally
$\varphi(0)=0$ by \eqref{eq:zero-full}; hence it is a rotation.
\end{proof}

\section{What remains}

The boundary extension is used twice: to turn fullness into the inner
boundary condition, and to prevent fiber mass from escaping through
infinitely many irregular sheets.  The proof does not extend to an arbitrary
analytic symbol by approximation.  Indeed, at a point where
$|\varphi'|=1$ for a hypothetical nonrotation fixing zero, its local term in
\eqref{eq:fiber} is strictly less than one, so other sheets are compulsory.
Those sheets may run to the boundary with nonvanishing hyperbolic derivative.

A later dissertation and article claim the unrestricted theorem
\cite{ShabazzThesis,ST}.  The available dissertation proof suppresses the
other sheets after choosing a local univalence disk.  A companion proof-gap
packet records that issue; it is not a counterexample to the theorem.  Thus
Theorem~\ref{thm:main} should be read as a rigorous partial result, not as a
replacement claim for the unrestricted classification.

\begin{thebibliography}{9}
\bibitem{AHP}
R. F. Allen, K. C. Heller, and M. A. Pons,
\emph{Isometric composition operators on the analytic Besov spaces},
J. Math. Anal. Appl. 414 (2014), 414--423; arXiv:2207.12634.

\bibitem{ShabazzThesis}
M. A. Shabazz,
\emph{Isometries of Besov Type Spaces among Composition Operators},
Ph.D. dissertation, University of Arkansas, 2015.

\bibitem{ST}
M. Shabazz and M. Tjani,
\emph{Isometries among composition operators on Besov type spaces},
Complex Variables and Elliptic Equations 61 (2016), 761--771,
doi:10.1080/17476933.2015.1118471.
\end{thebibliography}

\end{document}
